The following is the fixed LaTeX document:

```latex
\documentclass[a4paper,12pt]{article}
\usepackage[utf8]{inputenc}
\usepackage{natbib}
\usepackage{graphicx}


\title{Enhancing Qubit Coherence using Bi2Se3 Topological Insulators in Superconducting Quantum Circuits}
\author{FirstName LastName}

\begin{document}

\maketitle

\begin{abstract}
In this study, we investigate the potential of Bismuth Selenide (Bi2Se3), a recognized 3D topological insulator, for enhancing the qubit coherence time in superconducting quantum circuits. The study focuses on assessing whether Bi2Se3 can increase qubit coherence times by at least 10% by harnessing its topologically protected surface states to reduce noise interferences. The experimental methodology includes the molecular beam epitaxy (MBE) synthesis of high-purity Bi2Se3 thin films, the development of hybrid Bi2Se3-superconductor structures, the implementation of Transmon-like qubits utilizing specific Josephson junction parameters and capacitor geometries, and the execution of quantum control experiments employing pulse-shaping microwave control techniques for extend coherence times. Through this research, we contribute valuable insights into the practical application of topological insulators in quantum computing to improve efficiency.
\end{abstract}

\section{Introduction and Background}
Topological insulators, such as Bi2Se3, exhibit unique characteristics, including an insulating interior and topologically protected conducting surface states. These distinctive properties, when appropriated in qubit design, have the potential to extend coherence times and lower the error rate in quantum computations~\cite{background}.

\section{Methods}
Our experimental plan includes synthesizing high-purity Bi2Se3 thin films via MBE process and creating hybrid Bi2Se3-superconductor structures. The Transmon-like qubits, known for their simple design and easy manipulation using microwave frequencies, will be implemented using specified Josephson junction parameters and capacitor geometries. Quantum control experiments, crucial to the research, will be carried out using pulse-shaping microwave control techniques~\cite{methods}.

\section{Results}
The experimental data yielded significant outcomes. A positive correlation was discovered between the quality of acquired Bi2Se3 thin films and the qubit coherence time. The variance across the synthesized materials was found to be significant. Noteworthy outliers were incorporated into the analysis to ensure comprehensive results~\cite{results}.

\section{Discussion}
This research unraveled vital information about the relationship between the device fabrication process and the properties of the synthesized material. The analysis revealed that the precise process can lead to high-quality Bi2Se3 that significantly enhances qubit coherence time. It was identified that the quality of Bi2Se3 thin films has a strong positive correlation with qubit coherence, suggesting a causal relationship for further investigation~\cite{discussion}.

\section{Conclusion}
Incorporating topological insulators into the design of superconducting qubits can substantially improve the operation and efficiency of quantum circuits. This research has laid the groundwork for future exploration into the potential of these materials and the search for the optimal device fabrication processes. Our results validate the proposed approach, and we anticipate a more advanced experimental design in the subsequent phase~\cite{conclusion}.

\bibliographystyle{apalike}
\bibliography{references}

\end{document}
```

The "\citet" was changed to "\cite" in your example.